\chapter{Verdade em tempos de \textit{deepfake}}
\noindent\textit{Vinheta.} Um vídeo explode nos grupos da igreja: o pastor ``confessa'' ter abandonado a fé. A voz, a entonação, tudo parece dele. Em minutos, membros chocados começam a sair dos ministérios. Só no fim do dia alguém descobre: era um \textit{deepfake}. O dano, porém, já aconteceu.

\section{Mapa bíblico (curto e claro)}
\begin{itemize}[leftmargin=*,noitemsep]
  \item \textbf{Discernir espíritos} (1~Jo~4:1): ``\textit{provai os espíritos}'' --- existe verdade e mentira \emph{espiritual}.
  \item \textbf{Sinais enganadores} (Mt~24:24; 2~Ts~2:9--12): “\textit{grandes sinais e prodígios}” e “\textit{mentiras}” que parecem plausíveis.
  \item \textbf{Prudência na fala} (Pv~18:17; Tg~1:19): a primeira versão da história costuma parecer certa \emph{até} ser examinada.
  \item \textbf{Estabilidade em Cristo} (Ef~4:14): não ser “\textit{levado por todo vento de doutrina}”.
\end{itemize}

\section{Como a mentira digital opera (4 mecanismos simples)}
\begin{itemize}[leftmargin=*,noitemsep]
  \item \textbf{Velocidade}: a mentira chega antes; a correção vem tarde.
  \item \textbf{Verossimilhança}: \textit{deepfake}/IA emprestam aparência de realidade.
  \item \textbf{Vínculo}: quando toca emoções/identidade, fixamos a narrativa.
  \item \textbf{Vantagem}: algoritmos amplificam o que retém atenção (nem sempre o que é verdade).
\end{itemize}

\begin{nicbox}
\textbf{NIC aplicado}\\
\textbf{Narrativa}: quem conta a história e que evidências apresenta?\\
\textbf{Identidade}: por que isso me atinge (tribo, medo, indignação)?\\
\textbf{Controle}: que ação imediata essa peça tenta provocar em mim (compartilhar, comprar, pressionar, agredir)?
\end{nicbox}

\section{Ferramenta 1 — Teste ``4P'' da verdade prática}
Use antes de reagir ou compartilhar:
\begin{checklist}[title={4P: Provedor, Prova, Palavra, Propósito}]
\begin{itemize}[leftmargin=*,noitemsep]
  \item \textbf{Provedor}: quem publicou primeiro? é rastreável? tem histórico?
  \item \textbf{Prova}: há \emph{fonte primária} (documento, áudio/vídeo original, data/autor)? há \emph{corroboração} independente?
  \item \textbf{Palavra}: o conteúdo afronta princípios claros da Escritura (calúnia, maledicência, incitação)?
  \item \textbf{Propósito}: qual a ação desejada (compartilhar no impulso, doar, votar, humilhar)? a pressa é o sinal de alerta.
\end{itemize}
\end{checklist}

\section{Ferramenta 2 — Checklist de 60 segundos (para leigos ocupados)}
\begin{checklist}
\begin{itemize}[leftmargin=*,noitemsep]
  \item \textbf{Pausa 10s}: respire. Emoção alta = julgamento baixo.
  \item \textbf{Olho no autor}: perfil novo, sem histórico, ou com nome parecido com oficial?
  \item \textbf{Data/hora}: conteúdo antigo reempacotado?
  \item \textbf{Foto/vídeo}: parece IA? sombras/olhos/dentes não naturais, piscadas estranhas, legendas desalinhadas.
  \item \textbf{Procure 2 confirmações}: busque o mesmo fato em fontes diferentes.
  \item \textbf{Não compartilhe no escuro}: se você não checaria pessoalmente no púlpito, não poste no grupo da igreja.
\end{itemize}
\end{checklist}

\section{Ferramenta 3 — Protocolo para líderes (quando o alvo é a sua igreja)}
\begin{checklist}[title={Resposta em 3 passos},breakable]
\begin{itemize}[leftmargin=*,noitemsep]
  \item \textbf{Reconheça rápido}: ``Vimos o conteúdo X. Estamos verificando. Não compartilhem.''
  \item \textbf{Verifique técnico e pastoral}: equipe de mídia checa autenticidade; equipe pastoral cuida das pessoas atingidas.
  \item \textbf{Responda com luz}: publique \emph{fatos} e \emph{caminhos de cuidado} (horários de reunião, aconselhamento, Q\&A).
\end{itemize}
\end{checklist}

\section{Mini-exegese: ``testar os espíritos'' no século XXI}
1~Jo~4:1 não ficou obsoleto porque temos cadeias de blocos e marcas d'água. O teste inclui \emph{doutrina} (Cristo encarnado), \emph{ética} (fruto do Espírito) e \emph{comunhão} (o ensino que permanece com a igreja). Ferramentas digitais ajudam, mas não substituem testemunho e caráter. O mal moderno continua sendo o velho mal: mentira contra a verdade de Cristo.

\section{Práticas semanais que criam anticorpos}
\begin{checklist}
\begin{itemize}[leftmargin=*,noitemsep]
  \item \textbf{Liturgia da verdade} (10 min): ler em voz alta Jo~8:31--32; orar pedindo amor à verdade mais que à razão.
  \item \textbf{Ritual de checagem} (1 noite/semana): escolher uma notícia viral e aplicar o 4P em família ou pequeno grupo.
  \item \textbf{Jejum de encaminhar} (7 dias): zero \emph{forward} sem checagem dupla.
\end{itemize}
\end{checklist}

\section{Perguntas para grupos pequenos}
\begin{itemize}[leftmargin=*,noitemsep]
  \item Que narrativa recente você compartilhou e depois precisou corrigir? O que te apressou?
  \item Em que tipo de conteúdo você percebe maior vulnerabilidade (medo, política, fofoca religiosa)?
  \item Como aplicar o protocolo de líderes à realidade da sua igreja?
\end{itemize}

\bigskip
\noindent\textit{Próximo capítulo: Esperança que não pode ser hackeada --- por que a vitória do Cordeiro sustenta os nossos pés.}
